\input{preamble.tex}

\title{Лабораторная работа № 10. \\ Настройка прав доступа пользователей к ресурсам сети}
\author{Сущенко Алина Николаевна \\ НПИбд-01-23}

\institute{Российский университет дружбы народов имени Патриса Лумумбы}
\date{2025}

\begin{document}

\frame{\titlepage}

\begin{frame}
\frametitle{Цель работы}
\begin{itemize}
    \item Освоить настройку прав доступа пользователей к ресурсам сети с использованием стандартных и расширенных списков контроля доступа (Access Control Lists, ACL) в Cisco Packet Tracer.
\end{itemize}
\end{frame}

\begin{frame}
\frametitle{Топология сети и подключение администратора}
\begin{itemize}
    \item В топологию добавлен ноутбук \texttt{admin-ansusenko}, подключенный к порту 24 коммутатора \texttt{msk-ansusenko-sw-4}.
    \item Администратору назначен статический IP-адрес: 10.128.6.200, шлюз 10.128.6.1, DNS 10.128.0.5.
\end{itemize}
\centering
\includegraphics[width=0.7\textwidth]{../images/01.png}
\captionof{figure}{Обновленная топология сети с администратором}
\end{frame}

\begin{frame}
\frametitle{Настройки сетевого интерфейса администратора}
\begin{itemize}
    \item IPv4-адрес: 10.128.6.200
    \item Маска подсети: 255.255.255.0
    \item Шлюз по умолчанию: 10.128.6.1
    \item DNS-сервер: 10.128.0.5
\end{itemize}
\begin{center}
    \includegraphics[width=0.8\textwidth]{../images/02.png}
    \captionof{figure}{Настройки интерфейса FastEthernet0 администратора}
\end{center}
\end{frame}

\begin{frame}
\frametitle{Создание ACL servers-out}
\begin{itemize}
    \item Создан расширенный список доступа \texttt{servers-out} для управления доступом к серверам.
    \item Применен на интерфейсе \texttt{f0/0.3} маршрутизатора \texttt{msk-ansusenko-gw-1}.
    \item Первое правило: разрешить доступ всем пользователям к web-серверу (10.128.0.2) по порту 80.
\end{itemize}
\centering
\includegraphics[width=0.8\textwidth]{../images/03.png}
\captionof{figure}{Создание ACL servers-out и применение на интерфейсе}
\end{frame}

\begin{frame}
\frametitle{Добавление правил доступа для администратора}
\begin{itemize}
    \item Добавлены правила для доступа администратора (10.128.6.200) по FTP (порт 21) и Telnet (порт 23) к web-серверу.
\end{itemize}
\centering
\includegraphics[width=0.8\textwidth]{../images/04.png}
\captionof{figure}{Добавление правил доступа для администратора}
\end{frame}

\begin{frame}
\frametitle{Проверка FTP доступа администратора}
\begin{itemize}
    \item С устройства администратора выполнена команда \texttt{ftp 10.128.0.2}.
    \item FTP-доступ к web-серверу успешно работает.
\end{itemize}
\centering
\includegraphics[width=0.8\textwidth]{../images/05.png}
\captionof{figure}{Успешное подключение к FTP-серверу}
\end{frame}

\begin{frame}
\frametitle{Проверка FTP доступа пользователя}
\begin{itemize}
    \item С устройства \texttt{dep-ansusenko-1} выполнена попытка подключения по FTP.
    \item Доступ запрещён, что соответствует заданию --- FTP разрешён только администратору.
\end{itemize}
\centering
\includegraphics[width=0.8\textwidth]{../images/06.png}
\captionof{figure}{Ошибка подключения к FTP с устройства dep-ansusenko-1}
\end{frame}

\begin{frame}
\frametitle{Настройка доступа к файловому серверу}
\begin{itemize}
    \item Добавлены правила для доступа к файловому серверу (10.128.0.3):
    \begin{itemize}
        \item Внутренней сети (10.128.0.0/24) разрешён доступ по SMB (порт 445).
        \item Всем пользователям разрешён доступ по FTP (порт 21).
    \end{itemize}
\end{itemize}
\centering
\includegraphics[width=0.8\textwidth]{../images/07.png}
\captionof{figure}{Настройка доступа к файловому серверу}
\end{frame}

\begin{frame}
\frametitle{Настройка доступа к почтовому серверу}
\begin{itemize}
    \item Добавлены правила для доступа к почтовому серверу (10.128.0.4):
    \begin{itemize}
        \item Разрешён доступ по SMTP (порт 25).
        \item Разрешён доступ по POP3 (порт 110).
    \end{itemize}
\end{itemize}
\centering
\includegraphics[width=0.8\textwidth]{../images/08.png}
\captionof{figure}{Настройка доступа к почтовому серверу}
\end{frame}

\begin{frame}
\frametitle{Проверка работы DNS-сервера}
\begin{itemize}
    \item Выполнена команда \texttt{nslookup www.ansusenko.rudn.ru}.
    \item Выполнен \texttt{ping} по имени.
    \item DNS-сервер работает корректно, доступ к web-серверу по имени успешен.
\end{itemize}
\centering
\includegraphics[width=0.8\textwidth]{../images/09.png}
\captionof{figure}{Проверка работы DNS-сервера}
\end{frame}

\begin{frame}
\frametitle{Добавление правила для ICMP}
\begin{itemize}
    \item Добавлено правило разрешения ICMP в начало списка:
\end{itemize}
\ttfamily msk-ansusenko-gw-1(config-ext-nacl)\#1 permit icmp any any
\begin{itemize}
    \item Правило добавлено в начало списка, что демонстрирует управление порядком применения правил.
\end{itemize}
\centering
\includegraphics[width=0.8\textwidth]{../images/10.png}
\captionof{figure}{Добавление правила разрешения ICMP}
\end{frame}

\begin{frame}
\frametitle{Просмотр настроенных списков доступа}
\begin{itemize}
    \item Выполнена проверка всех настроенных списков доступа.
\end{itemize}
\centering
\includegraphics[width=0.8\textwidth]{../images/11.png}
\captionof{figure}{Просмотр настроенных списков доступа}
\end{frame}

\begin{frame}
\frametitle{Настройка доступа для сети Other}
\begin{itemize}
    \item Создан список \texttt{other-in} и применён на интерфейсе f0/0.104 на входящий трафик.
    \item Разрешены любые действия только для администратора (10.128.6.200).
\end{itemize}
\centering
\includegraphics[width=0.8\textwidth]{../images/12.png}
\captionof{figure}{Настройка доступа для сети Other}
\end{frame}

\begin{frame}
\frametitle{Настройка доступа к сети управления оборудованием}
\begin{itemize}
    \item Создан список \texttt{management-out} для доступа к сети управления оборудованием (10.128.1.0/24).
\end{itemize}
\centering
\includegraphics[width=0.8\textwidth]{../images/13.png}
\captionof{figure}{Настройка доступа к сети управления оборудованием}
\end{frame}

\begin{frame}
\frametitle{Проверка доступа не с администратора}
\begin{itemize}
    \item Выполнена проверка доступа к серверам с устройства \texttt{dep-ansusenko-1}.
    \item Доступ запрещён в соответствии с настроенными правилами.
\end{itemize}
\centering
\includegraphics[width=0.8\textwidth]{../images/14.png}
\captionof{figure}{Проверка доступа к серверам (не администратор)}
\end{frame}

\begin{frame}
\frametitle{Проверка доступа с администратора}
\begin{itemize}
    \item Выполнена проверка доступа администратора к серверам.
    \item Весь необходимый доступ разрешён.
\end{itemize}
\centering
\includegraphics[width=0.8\textwidth]{../images/15.png}
\captionof{figure}{Проверка доступа администратора к серверам}
\end{frame}

\begin{frame}
\frametitle{Добавление нового администратора в топологию}
\begin{itemize}
    \item В соответствии с заданием самостоятельной работы добавлен новый ноутбук в топологию сети Павловской.
\end{itemize}
\centering
\includegraphics[width=0.8\textwidth]{../images/18.png}
\captionof{figure}{Добавление нового ноутбука в топологию сети Павловской}
\end{frame}

\begin{frame}
\frametitle{Настройка прав доступа для нового администратора}
\begin{itemize}
    \item Настроены права доступа для нового администратора.
\end{itemize}
\centering
\includegraphics[width=0.8\textwidth]{../images/16.png}
\captionof{figure}{Права доступа для нового администратора}
\end{frame}

\begin{frame}
\frametitle{Проверка настроек для нового администратора}
\begin{itemize}
    \item Выполнена проверка списков и настроек для нового администратора.
\end{itemize}
\centering
\includegraphics[width=0.8\textwidth]{../images/17.png}
\captionof{figure}{Проверка доступа администратора к серверам}
\end{frame}

\begin{frame}
\frametitle{Настройки интерфейса второго администратора}
\begin{itemize}
    \item IPv4-адрес: 10.128.6.201
    \item Маска подсети: 255.255.255.0
    \item Шлюз по умолчанию: 10.128.6.1
    \item DNS-сервер: 10.128.0.5
\end{itemize}
\centering
\includegraphics[width=0.8\textwidth]{../images/19.png}
\captionof{figure}{Настройки администратора на Павловской}
\end{frame}

\begin{frame}
\frametitle{Добавление правил для второго администратора}
\begin{itemize}
    \item Добавлены правила для администратора 10.128.6.201 в списки \texttt{other-in} и \texttt{management-out}.
\end{itemize}
\centering
\includegraphics[width=0.8\textwidth]{../images/20.png}
\captionof{figure}{Добавление правил для администратора на Павловской}
\end{frame}

\begin{frame}
\frametitle{Итоговый вывод ACL}
\begin{itemize}
    \item Выполнен вывод \texttt{show access-lists} с правилами для обоих администраторов (10.128.6.200 и 10.128.6.201).
\end{itemize}
\centering
\includegraphics[width=0.8\textwidth]{../images/21.png}
\captionof{figure}{Итоговый вывод ACL после настройки второго администратора}
\end{frame}

\begin{frame}
\frametitle{Проверка доступа второго администратора}
С устройства администратора 10.128.6.201 выполнены:
\begin{itemize}
    \item \texttt{ping 10.128.6.1} --- успешен
    \item \texttt{ping 10.128.0.2} --- 75\% успешных ответов
    \item \texttt{ftp 10.128.0.2} --- успешен
    \item \texttt{ping 10.128.1.1} --- успешен
\end{itemize}
\textbf{Вывод:} Второй администратор имеет полный доступ к ресурсам сети.
\centering
\includegraphics[width=0.7\textwidth]{../images/22.png}
\captionof{figure}{Проверка доступа администратора на Павловской}
\end{frame}

\begin{frame}
\frametitle{Вопрос 1}
\parbox{\textwidth}{
\textbf{Вопрос:} Как задать действие правила для конкретного протокола?

\textbf{Ответ:} Для задания действия правила для конкретного протокола в расширенном ACL используется следующий синтаксис:

\ttfamily access-list <номер\_или\_имя> <permit|deny> <протокол> <источник> <порт\_источника> <назначение> <порт\_назначения>

\textbf{Пример из работы:}

\ttfamily permit tcp any host 10.128.0.2 eq 80

Данное правило разрешает (permit) протокол TCP любому источнику (any) к хосту 10.128.0.2 на порт 80 (eq 80).

\textbf{Возможные протоколы:} tcp, udp, icmp, ip, gre и другие.
}
\end{frame}

\begin{frame}
\frametitle{Вопрос 2}
\parbox{\textwidth}{
\textbf{Вопрос:} Как задать действие правила сразу для нескольких портов?

\textbf{Ответ:} Существует несколько способов задать действие правила сразу для нескольких портов:

\textbf{1. Использование оператора диапазона range:}

\ttfamily permit tcp any host 10.128.0.4 range 25 110

Разрешает доступ к портам с 25 по 110.

\textbf{2. Использование нескольких правил:}

\ttfamily permit tcp any host 10.128.0.4 eq 25

\ttfamily permit tcp any host 10.128.0.4 eq 110

\textbf{3. Использование операторов lt (меньше), gt (больше), neq (не равно):}

\ttfamily permit tcp any host 10.128.0.4 lt 1024

Разрешает доступ к портам меньше 1024 (привилегированные порты).
}
\end{frame}

\begin{frame}
\frametitle{Вопрос 3}
\parbox{\textwidth}{
\textbf{Вопрос:} Как узнать номер правила в списке прав доступа?

\textbf{Ответ:} Для просмотра номеров правил в списке доступа используется команда \texttt{show access-lists} в привилегированном режиме (enable).

\textbf{Пример вывода из работы:}

\ttfamily msk-ansusenko-gw-1\# show access-lists

\ttfamily Extended IP access list servers-out

\ttfamily \ \ \ 10 permit tcp any host 10.128.0.2 eq 80

\ttfamily \ \ \ 20 permit tcp host 10.128.6.200 host 10.128.0.2 eq 21

\ttfamily \ \ \ 30 permit tcp host 10.128.6.200 host 10.128.0.2 eq 23

\textbf{Особенности:}
\begin{itemize}
    \item Номера правил отображаются в первом столбце (10, 20, 30).
    \item Если номера не заданы явно, Cisco присваивает их автоматически с шагом 10.
    \item Номера правил позволяют управлять порядком применения (меньший номер --- выше приоритет).
\end{itemize}
}
\end{frame}

\begin{frame}
\frametitle{Вопрос 4}
\parbox{\textwidth}{
\textbf{Вопрос:} Каким образом можно изменить порядок применения правил в списке контроля доступа?

\textbf{Ответ:} Существует несколько способов изменения порядка применения правил:

\textbf{1. Добавление правила с конкретным номером:}

\ttfamily ip access-list extended servers-out

\ttfamily 1 permit icmp any any

Правило с номером 1 будет первым, независимо от существующих правил.

\textbf{2. Удаление и повторное добавление правил в нужном порядке:}

\ttfamily no access-list extended servers-out

\ttfamily access-list extended servers-out

\ttfamily permit icmp any any

\ttfamily permit tcp any host 10.128.0.2 eq 80

\textbf{3. Изменение порядка с помощью нумерации (в именованных ACL):}

\ttfamily ip access-list extended servers-out

\ttfamily 15 permit tcp host 10.128.6.200 host 10.128.0.2 eq 21

\ttfamily no 20
}
\end{frame}

\begin{frame}
\frametitle{Выводы}
\begin{itemize}
    \item В результате выполнения лабораторной работы приобрели практические навыки по настройке прав доступа пользователей к ресурсам сети с использованием расширенных списков контроля доступа Cisco.
    \item Были настроены ACL \texttt{servers-out}, \texttt{other-in} и \texttt{management-out}, обеспечивающие требуемые политики доступа к web-серверу, файловому серверу, почтовому серверу, DNS-серверу, а также ограничения для сети Other и доступ к сети управления оборудованием.
\end{itemize}
\end{frame}

\end{document}
